\documentclass[12pt,a4paper]{article}
\usepackage[utf8]{inputenc}
\usepackage[T5]{fontenc}
\usepackage{amsmath, amssymb}
\usepackage{graphicx}
\usepackage{ifthen}

\newboolean{showsolutions}
\setboolean{showsolutions}{true} % Đổi thành false để ẩn đáp án khi in PDF cho học sinh

% --- ĐỊNH NGHĨA CÁC LỆNH ẢO CHO CÔNG CỤ LATEX2JSON ---
\newcommand{\doctitle}[1]{\begin{center}\Large\textbf{#1}\end{center}}
\newcommand{\docdesc}[1]{\textbf{Mô tả:} #1\par}
\newcommand{\docgrade}[1]{\textbf{Lớp:} #1\par}
\newcommand{\docstatus}[1]{\textbf{Trạng thái:} #1\par}
\newcommand{\doctype}[1]{\textbf{Loại:} #1\par}
\newcommand{\doctopics}[1]{\textbf{Chủ đề:} #1\par}

\newenvironment{textblock}{\vspace{1em}}{\par}

\newenvironment{lesson}[2]{
	\vspace{1em}\noindent\textbf{\Large #1}\\
	\textit{#2}\par\vspace{0.5em}
}{}

\newcommand{\image}[3]{
	\begin{center}
		\includegraphics[width=#1\textwidth]{#3} \\
		\textit{#2}
	\end{center}
}

\newenvironment{quiz}[1]{
	\vspace{1em}\noindent\textbf{\Large Kiểm tra: #1}\par
}{}

\newenvironment{mcq}[4]{
	\vspace{1em}\noindent\textbf{Câu hỏi:} #1 \\
	\ifthenelse{\boolean{showsolutions}}{
		\textit{(Đáp án đúng: #2, Điểm: #3) - Giải thích: #4}
	}{}
	\begin{itemize}
	}{
	\end{itemize}
}
\newcommand{\option}[1]{\item #1}

\newenvironment{truefalse}[3]{
	\vspace{1em}\noindent\textbf{Đúng/Sai:} #1 \\
	\ifthenelse{\boolean{showsolutions}}{
		\textit{(Điểm: #2) - Giải thích: #3}
	}{}
	\begin{itemize}
	}{
	\end{itemize}
}
\newcommand{\statement}[2]{\item [\textbf{#1}] #2}

\newcommand{\shortanswer}[4]{
    \vspace{1em}\noindent\textbf{Trả lời ngắn (Điểm: #3):}

    \noindent #1

	\ifthenelse{\boolean{showsolutions}}{
		\vspace{0.5em}
		\noindent\textit{Đáp án: #2 - Giải thích:}
		\noindent #4\par
	}{}
}

\newenvironment{shortanswer}[4]{
    \vspace{1em}\noindent\textbf{Trả lời ngắn (Điểm: #3):}

    \noindent #1

	\ifthenelse{\boolean{showsolutions}}{
		\vspace{0.5em}
		\noindent\textit{Đáp án: #2 - Giải thích:}
		\noindent #4\par
	}{}
}{}

\newcommand{\essay}[3]{
    \vspace{1em}\noindent\textbf{Tự luận (Điểm: #2):}
    
    \noindent #1
    
	\ifthenelse{\boolean{showsolutions}}{
		\vspace{0.5em}
		\noindent\textbf{Gợi ý / Đáp án:}
		\noindent #3\par
	}{}
}

% =============================================================================

\begin{document}
	
	\doctitle{TÍNH ĐƠN ĐIỆU CỦA HÀM SỐ}
	\docdesc{Tóm tắt lý thuyết về tính đơn điệu của hàm số}
	\docgrade{Lớp 12}
	\docstatus{draft}
	\doctype{theory}
	\doctopics{tinh-don-dieu-cua-ham-so}
	
	\begin{lesson}{PHẦN I. TÓM TẮT LÝ THUYẾT}{}
		
		\textbf{1. Định nghĩa:} Cho hàm số $y = f(x)$ xác định trên $K$ ($K$ là khoảng, đoạn hoặc nửa khoảng).
		\begin{itemize}
			\item Hàm số $y = f(x)$ được gọi là \textbf{đồng biến (tăng)} trên $K$ nếu 
			$$\forall x_1, x_2 \in K, x_1 < x_2 \Rightarrow f(x_1) < f(x_2).$$
			\item Hàm số $y = f(x)$ được gọi là \textbf{nghịch biến (giảm)} trên $K$ nếu 
			$$\forall x_1, x_2 \in K, x_1 < x_2 \Rightarrow f(x_1) > f(x_2).$$
		\end{itemize}
		
		\textbf{Chú ý.}
		\begin{itemize}
			\item Nếu hàm số \textbf{đồng biến} trên $K$ thì đồ thị của hàm số \textbf{đi lên} từ trái sang phải (Hình 1).
			\item Nếu hàm số \textbf{nghịch biến} trên $K$ thì đồ thị của hàm số \textbf{đi xuống} từ trái sang phải (Hình 2).
			% ─── KHỐI HÌNH ẢNH (image block) ─────────────────────────────────────────────
			% Cú pháp: \image{tỉ lệ kích cỡ (VD: 0.5, 0.7...)}{caption}{đường dẫn file ảnh}
			\image{0.7}{Hình 1: Đồ thị hàm số đồng biến}{do_thi_don_dieu.png}
			\item Hàm số đồng biến hoặc nghịch biến trên $K$ gọi chung là \textbf{\textit{đơn điệu}} trên $K$. Việc tìm các khoảng đồng biến, nghịch biến của hàm số còn được gọi là tìm các khoảng đơn điệu (hay xét tính đơn điệu) của hàm số.
			\item Khi xét tính đơn điệu của hàm số mà không chỉ rõ tập $K$ thì ta hiểu là xét trên tập xác định của hàm số đó.
		\end{itemize}
		
		\textbf{2. Liên hệ giữa đạo hàm và tính đơn điệu:} Cho hàm số $y = f(x)$ có đạo hàm trên $K$.
		\begin{itemize}
			\item Nếu $f'(x) > 0, \forall x \in K$ thì hàm số $f(x)$ \textbf{đồng biến} trên $K$.
			\item Nếu $f'(x) < 0, \forall x \in K$ thì hàm số $f(x)$ \textbf{nghịch biến} trên $K$.
			\item Nếu $f'(x) = 0, \forall x \in K$ thì hàm số $f(x)$ \textbf{không đổi} trên $K$ hoặc $f(x)$ là \textbf{hàm hằng} trên $K$.
		\end{itemize}
	\end{lesson}
	
\end{document}
